\documentclass{beamer}
\usepackage{graphicx}
\graphicspath{{./figs/}}
\usepackage{amsmath, amssymb}
\usepackage{tcolorbox}
\usepackage{caption}
\usepackage{subcaption}
\usepackage{multirow}


\title{CE 401/5501: Applied Artificial Intelligence}
\subtitle{Topic 4: Simple yet Significant Logistic Regression}
\author{Dr. ZhiQiang Chen}
\date{}

\begin{document}

\begin{frame}
  \titlepage
\end{frame}

\begin{frame}{Outline}
  \tableofcontents
\end{frame}

\section{Introduction \& Historical Context}
\begin{frame}{Introduction}
  \begin{itemize}
    \item Understand logistic regression as a discriminative, supervised classifier.
    \item Generalize its formulation to comprehend other classifiers.
  \end{itemize}
\end{frame}

\begin{frame}{Historical Context \& Significance}
  \begin{tcolorbox}[colback=blue!5!white,colframe=blue!75!black]
    \textbf{Historical Context:}\\[1ex]
    Logistic regression originated as a statistical method for binary outcomes (e.g., survival vs.\ non-survival) and evolved into a widely used machine learning classifier.\\[1ex]
    \textbf{Significance:}
    \begin{itemize}
      \item \textbf{Simplicity \& Interpretability}
      \item \textbf{Computational Efficiency}
      \item \textbf{Probability Estimation}
    \end{itemize}
  \end{tcolorbox}
\end{frame}

\section{Sigmoid and Logistic Functions}
\begin{frame}{Sigmoid / Logistic Function}
  \begin{block}{Definition}
    The Sigmoid function is a special case of the logistic function:
    \[
    \mathrm{Sigmoid:}\ f(x) = \frac{1}{1 + \exp(-x)}
    \]
      \[
     \mathrm{Logistic:}\  f[u(x)] = L\ \frac{1}{1 + \exp[-u(x)]}
    \]
  \end{block}
  \begin{itemize}
    \item Continuous and monotonically increasing.
    \item $f(0)=\frac{1}{2}$, and limits:
      \[
      \lim_{x \to -\infty} f(x) = 0,\quad \lim_{x \to \infty} f(x) = 1.
      \]
  \end{itemize}
\end{frame}

\begin{frame}{Sigmoid Function: Derivative and Figure}
  \begin{block}{Derivative}
    \[
    \frac{df(x)}{dx} = f(x)[1-f(x)]
    \]
  \end{block}
  \begin{figure}
    \centering
    \includegraphics[width=0.5\textwidth]{figs/topic04fig01.png}
    \caption{Sigmoid function.}
  \end{figure}
\end{frame}

\begin{frame}{Logistic Function as an ML Model}
  \begin{itemize}
    \item Transform input via \( u(x) = w_0 + w_1 x \).
    \item Logistic model: 
      \[
      f(x,\mathbf{w}) = \frac{1}{1+\exp(-w_0 - w_1 x)}
      \]
    \item Decision making via threshold \( f(x) \ge \frac{1}{2} \) or \( u(x) \ge 0 \).
  \end{itemize}
\end{frame}

\section{Learning from Data and Decision Function}
\begin{frame}{Learning from Data}
  \begin{itemize}
    \item Given a dataset:
      \[
      \mathbf{D} = \{(x_k, c_k) \mid k=1,2,\ldots,N\}
      \]
    \item Toy example (soil-health experiment):
      \[
      \mathbf{D}_s = \{(3.0,0), (5.0,0), (7.0,1)\}
      \]
    \item Decision boundary defined in the original and transformed spaces.
  \end{itemize}
\end{frame}

\begin{frame}{Decision Function Variants}
  \begin{itemize}
    \item \textbf{Linear:} \( u(x) = w_0 + w_1 x \)
    \item \textbf{Nonlinear:} Polynomial transformation \( u(x)=w_0+w_1 x+\beta_2x^2+\cdots \)
    \item \textbf{Multidimensional:} For \(\mathbf{x} \in \mathbb{R}^d\):
      \[
      u(\mathbf{x}) = \langle \mathbf{w},\mathbf{x}\rangle + w_0 = \sum_{i=1}^d w_i x_i + w_0
      \]
    \item \textbf{Kernel Methods:} Nonlinear transformations in high-dimensional spaces.
  \end{itemize}
\end{frame}

\begin{frame}{Interpreting the Output Probability}
  \begin{itemize}
    \item The logistic function outputs a value \( f[u(x)] \) that can be directly interpreted as a probability.
    \item For example:
      \begin{itemize}
        \item \( f = 0.75 \) implies that the probability of \( m(x)=1 \) is 75\% (and \( m(x)=0 \) is 25\%).
        \item \( f = 0.25 \) implies that the probability of \( m(x)=1 \) is 25\% (or equivalently, \( m(x)=0 \) is 75\%).
      \end{itemize}
    \item This probabilistic interpretation allows us to use \( f[u(x)] \) directly as a classifier.
  \end{itemize}
\end{frame}

\begin{frame}{Learning the Structure of \( u(x) \) from Data}
  \begin{itemize}
    \item Since \( f[u(x)] \) is used directly as a classifier, the goal is to learn the underlying function \( u(x) \) from data.
    \item In the simplest case, \( u(x) \) is assumed to be linear:
      \[
      u(x) = w_1 x + w_0.
      \]
    \item Thus, the learning problem reduces to estimating the parameters \( w_1 \) and \( w_0 \) based on the training data.
  \end{itemize}
\end{frame}


\begin{frame}{Estimating the Model Parameters - a thought process}
  \begin{itemize}
    \item We assume a parametric structure for the decision function: 
      \[
      u(x) = w_1 x + w_0.
      \]
    \item Given a set of labeled training data \( \{(x_1, c_1), (x_2, c_2), \dots, (x_N, c_N)\} \), these data points constrain the possible values of \( w_1 \) and \( w_0 \).
    \item Among many possible parameter values, there exists an optimal set that best explains the observed labels.
    \item To find these optimal weights, we define an objective (loss) function that quantifies the discrepancy between the predicted outputs \( f[u(x)] \) and the true labels.
    \item The goal is to minimize this loss function to obtain the best estimates of \( w_1 \) and \( w_0 \).
  \end{itemize}
\end{frame}


\section{Loss Function and Optimization}
\begin{frame}{Loss Function: Log-Loss}
  \begin{block}{Loss for a single sample}
    \[
    l(f_k) = -c_k\ln(f_k) - (1-c_k)\ln(1-f_k)
    \]
  \end{block}
  \begin{itemize}
    \item Total loss over dataset:
      \[
      L = -\sum_k \left[c_k\ln(f_k) + (1-c_k)\ln(1-f_k)\right]
      \]
    \item Encourages high confidence in predictions.
  \end{itemize}
\end{frame}

\begin{frame}{Optimization: Gradient-Based Methods}
  \begin{itemize}
    \item Objective: Find \(\mathbf{w}^* = \arg\min_{\mathbf{w}} L\).
    \item Gradient for a single sample:
      \[
      \nabla_\mathbf{w} l(\mathbf{w}) = -(c-f)
      \begin{bmatrix} 1 \\ x \end{bmatrix}
      \]
    \item Total gradient:
      \[
      \nabla_\mathbf{w} L =
      \begin{bmatrix}
      \sum_k -(c_k - f_k) \\
      \sum_k -(c_k - f_k)x_k
      \end{bmatrix}
      \]
    \item Use iterative methods (e.g., gradient descent or Newton's method) to update \(\mathbf{w}\).
  \end{itemize}
\end{frame}

\section{Toy Example}
\begin{frame}{Toy Example: Soil-Health Dataset}
  \begin{itemize}
    \item Dataset: 
      \[
      \mathcal{D} = \{(0,3.0), (0,5.0), (1,7.0)\}
      \]
    \item Logistic model:
      \[
      f(x,\mathbf{w}) = \frac{1}{1+\exp(-w_0-w_1x)}
      \]
    \item Total loss:
      \[
      L = -\ln\left[(1-f_1)(1-f_2)f_3\right]
      \]
      where
      \[
      f_1 = \frac{1}{1+\exp(-w_0-3w_1)},\quad
      f_2 = \frac{1}{1+\exp(-w_0-5w_1)}
      \]
      \[
      f_3 = \frac{1}{1+\exp(-w_0-7w_1)}.
      \]
  \end{itemize}
\end{frame}

\begin{frame}{Toy Example: Optimization and Decision Boundary}
  \begin{itemize}
    \item Parameter estimates (from numerical optimization):
      \[
      w_0 = -5.0718,\quad w_1 = 0.8186.
      \]
    \item Decision boundary: 
      \[
      0.8186\,x - 5.0718 = 0 \quad \Rightarrow \quad x^* \approx 6.1961.
      \]
    \item For \(x > x^*\), predict class 1; for \(x < x^*\), predict class 0.
  \end{itemize}
  \begin{figure}
    \centering
    \includegraphics[width=0.5\textwidth]{figs/topic04fig07.png}
    \caption{Loss function surface in terms of \(w_0\) and \(w_1\).}
  \end{figure}
\end{frame}

\section{Conclusion}
\begin{frame}{Conclusion}
  \begin{itemize}
    \item \textbf{Recap:} We covered the theory and practical implementation of logistic regression.
    \item \textbf{Key Points:}
      \begin{itemize}
        \item Sigmoid and logistic functions provide a probabilistic framework.
        \item Decision functions transform input data for effective classification.
        \item The log-loss function guides model training via gradient-based optimization.
      \end{itemize}
    \item \textbf{Practical Insight:} The toy example illustrated how parameter estimation yields a meaningful decision boundary.
    \item \textbf{Overall:} Logistic regression remains a fundamental, interpretable model in both statistics and modern AI.
  \end{itemize}
\end{frame}

\end{document}
